\documentclass[11pt,a4paper]{article}
\usepackage[utf8]{inputenc}
\usepackage[french]{babel}
\usepackage{amsmath,amssymb,amsthm}
\usepackage{geometry}
\geometry{hmargin=2.5cm,vmargin=3cm}
\usepackage{tikz}
\usetikzlibrary{cd,positioning,shapes}
\usepackage{listings}
\usepackage{xcolor}

\lstset{
    language=Python,
    basicstyle=\ttfamily\small,
    keywordstyle=\color{blue}\bfseries,
    stringstyle=\color{red},
    commentstyle=\color{gray}\itshape,
    numbers=left,
    numberstyle=\tiny\color{gray},
    breaklines=true,
    frame=single,
    showstringspaces=false,
    literate={é}{{\'e}}1 {è}{{\`e}}1 {à}{{\`a}}1 {ê}{{\^e}}1 {î}{{\^i}}1 {ç}{{\c{c}}}1 {ï}{{\"i}}1 {É}{{\'E}}1
}

\title{\Huge ALGÈBRES ET TRIBUS (SIGMA-ALGÈBRES) \\ \large Cours magistral, exercices corrigés et travaux pratiques}
\author{\Large Charles EDOU NZE}
\date{\today}

\begin{document}
\maketitle
\tableofcontents
\newpage

\part{Cours Magistral}

\section{Genèse et motivation conceptuelle}

La théorie de la mesure vise à assigner un "volume" (ou longueur, aire, masse, probabilité) aux sous-ensembles d'un espace donné. Intuitivement, on souhaiterait pouvoir mesurer n'importe quel sous-ensemble de $\mathbb{R}$. Cependant, l'axiome du choix et la géométrie de l'espace tridimensionnel conduisent à des paradoxes majeurs, tel que le paradoxe de Banach-Tarski, démontrant qu'il est impossible de définir une mesure consistante et invariante par translation sur toutes les parties de $\mathbb{R}^n$.

Pour pallier ce problème fondamental, l'approche axiomatique moderne restreint la collection des sous-ensembles que l'on s'autorise à mesurer. On ne mesure pas tout, mais seulement une famille d'ensembles dits "mesurables". Cette famille doit posséder une structure algébrique robuste vis-à-vis des opérations ensemblistes fondamentales : le passage au complémentaire et l'union dénombrable, permettant ainsi de définir des processus limites continus. C'est l'essence même du concept de tribu, ou $\sigma$-algèbre, introduit formellement par Émile Borel et Henri Lebesgue.

\section{Algèbres et Tribus}

Soit $X$ un ensemble non vide.

\subsection{Définitions et Théorèmes}

\textbf{Définition (Algèbre) :}
Une famille $\mathcal{A}$ de parties de $X$ est une \textbf{algèbre} sur $X$ si elle vérifie les trois propriétés suivantes :
1. $X \in \mathcal{A}$.
2. Stabilité par passage au complémentaire : Si $A \in \mathcal{A}$, alors $X \setminus A \in \mathcal{A}$.
3. Stabilité par union finie : Si $A, B \in \mathcal{A}$, alors $A \cup B \in \mathcal{A}$.

\textbf{Définition (Tribu ou $\sigma$-algèbre) :}
Une famille $\mathcal{F}$ de parties de $X$ est une \textbf{tribu} (ou $\sigma$-algèbre) sur $X$ si elle vérifie :
1. $X \in \mathcal{F}$.
2. Stabilité par passage au complémentaire : Si $A \in \mathcal{F}$, alors $X \setminus A \in \mathcal{F}$.
3. Stabilité par union dénombrable : Si $(A_n)_{n \in \mathbb{N}}$ est une suite d'éléments de $\mathcal{F}$, alors $\bigcup_{n \in \mathbb{N}} A_n \in \mathcal{F}$.

Le couple $(X, \mathcal{F})$ est alors appelé un \textbf{espace mesurable}.

\textbf{Exemples concrets (Tribus) :}
1. La tribu triviale $\mathcal{F} = \{\emptyset, X\}$ est la plus petite tribu sur $X$.
2. La tribu discrète $\mathcal{F} = \mathcal{P}(X)$ (l'ensemble des parties de $X$) est la plus grande tribu sur $X$.
3. Soit $A \subset X$. La famille $\mathcal{F} = \{\emptyset, A, X \setminus A, X\}$ est une tribu.
4. Sur $X = \{1, 2, 3\}$, la famille $\{\emptyset, \{1\}, \{2, 3\}, X\}$ est une tribu.
5. Soit $X$ un ensemble infini. La famille des parties $A \subset X$ telles que $A$ est finie ou $X \setminus A$ est finie forme une algèbre, mais en général pas une tribu.
6. La famille des parties $A \subset X$ telles que $A$ est dénombrable ou $X \setminus A$ est dénombrable forme une tribu.

\textbf{Exemples concrets (Algèbre non Tribu) :}
7. Sur $X = \mathbb{R}$, l'ensemble des réunions finies d'intervalles de la forme $]a, b]$, $]-\infty, a]$, ou $]a, +\infty[$ est une algèbre, mais non une tribu (l'union dénombrable de $]0, 1-1/n]$ donne $]0, 1[$ qui n'est pas de cette forme de base si l'on se restreint à des intervalles semi-ouverts de manière stricte, bien que l'algèbre soit définie différemment usuellement, disons l'algèbre engendrée par les intervalles semi-ouverts). Un meilleur exemple : l'algèbre des parties finies ou cofinies de $\mathbb{N}$. L'union des singletons $\{2n\}$ est l'ensemble des pairs, qui n'est ni fini ni cofini.

\begin{tikzpicture}
\node[draw, circle, minimum size=3cm] (X) at (0,0) {};
\node at (0, 1) {$X$};
\draw[fill=blue!20] (-0.5,-0.5) circle (0.8);
\node at (-0.5,-0.5) {$A$};
\draw[fill=red!20] (0.8,-0.2) circle (0.6);
\node at (0.8,-0.2) {$B$};
\node at (0, -2) {Dans une tribu, $A \cup B, A \cap B, X \setminus A$ sont inclus.};
\end{tikzpicture}

\subsection{Démonstrations et Propriétés}

\textbf{Propriété :} Une tribu est stable par intersection dénombrable.
\textit{Preuve détaillée :} Soit $(A_n)_{n \in \mathbb{N}}$ une suite d'éléments de $\mathcal{F}$. Nous voulons montrer que $\bigcap_{n \in \mathbb{N}} A_n \in \mathcal{F}$.
D'après les lois de De Morgan, le complémentaire d'une intersection est l'union des complémentaires :
$$ X \setminus \left( \bigcap_{n \in \mathbb{N}} A_n \right) = \bigcup_{n \in \mathbb{N}} (X \setminus A_n) $$
Puisque $\mathcal{F}$ est une tribu :
- Pour tout $n$, $A_n \in \mathcal{F}$, donc par stabilité par passage au complémentaire, $X \setminus A_n \in \mathcal{F}$.
- Par stabilité par union dénombrable, $\bigcup_{n \in \mathbb{N}} (X \setminus A_n) \in \mathcal{F}$.
- Enfin, par stabilité par passage au complémentaire appliquée à cette union, le complémentaire du complémentaire, c'est-à-dire $\bigcap_{n \in \mathbb{N}} A_n$, appartient à $\mathcal{F}$.

\textbf{Propriété :} Toute tribu est une algèbre.
\textit{Preuve détaillée :} La stabilité par union finie s'obtient en complétant une suite finie $A_1, \dots, A_k$ par l'ensemble vide : pour $n > k$, posons $A_n = \emptyset$. Puisque $\emptyset = X \setminus X \in \mathcal{F}$, la suite $(A_n)$ est dans $\mathcal{F}$, et $\bigcup_{n \in \mathbb{N}} A_n = \bigcup_{i=1}^k A_i \in \mathcal{F}$.

\section{Tribu engendrée et Tribu de Borel}

\subsection{Définitions et Théorèmes}

Il est souvent difficile de décrire explicitement tous les éléments d'une tribu. On la définit plutôt à partir d'une famille d'ensembles générateurs.

\textbf{Définition (Tribu engendrée) :}
Soit $\mathcal{C}$ une classe (famille) de parties de $X$. La \textbf{tribu engendrée} par $\mathcal{C}$, notée $\sigma(\mathcal{C})$, est la plus petite tribu contenant $\mathcal{C}$. Formellement, c'est l'intersection de toutes les tribus contenant $\mathcal{C}$ :
$$ \sigma(\mathcal{C}) = \bigcap \{ \mathcal{F} \text{ tribu sur } X \mid \mathcal{C} \subset \mathcal{F} \} $$

\textbf{Définition (Tribu borélienne) :}
Soit $(X, \mathcal{T})$ un espace topologique ($\mathcal{T}$ étant la famille des ouverts). La \textbf{tribu borélienne}, ou tribu de Borel, notée $\mathcal{B}(X)$, est la tribu engendrée par les ouverts de $X$ : $\mathcal{B}(X) = \sigma(\mathcal{T})$.
Les éléments de $\mathcal{B}(X)$ sont appelés les \textbf{boréliens}.

\textbf{Exemples concrets (Tribus engendrées) :}
1. Si $\mathcal{C} = \{A\}$ avec $A \subsetneq X, A \neq \emptyset$, alors $\sigma(\mathcal{C}) = \{\emptyset, A, X \setminus A, X\}$.
2. Sur $\mathbb{R}$, $\mathcal{B}(\mathbb{R})$ est générée par les ouverts $]a, b[$.
3. L'ensemble $\{a\}$ est un borélien car $\{a\} = \bigcap_{n \in \mathbb{N}^*} ]a - \frac{1}{n}, a + \frac{1}{n}[$.
4. Tout fermé est un borélien car complémentaire d'un ouvert.
5. Sur un ensemble fini $X$, si on partitionne $X$ en $\{A_i\}$, la tribu engendrée par cette partition contient exactement les unions de ces éléments $A_i$.

\subsection{Démonstrations}

\textbf{Théorème :} La tribu engendrée $\sigma(\mathcal{C})$ existe et est bien une tribu.
\textit{Preuve détaillée :}
1. L'ensemble des tribus contenant $\mathcal{C}$ n'est pas vide, car la tribu des parties $\mathcal{P}(X)$ contient $\mathcal{C}$ et est une tribu.
2. Montrons que l'intersection d'une famille quelconque de tribus $(\mathcal{F}_i)_{i \in I}$ est une tribu.
   - $\forall i \in I, X \in \mathcal{F}_i$, donc $X \in \bigcap_{i \in I} \mathcal{F}_i$.
   - Si $A \in \bigcap_{i \in I} \mathcal{F}_i$, alors $\forall i \in I, A \in \mathcal{F}_i$. Comme chaque $\mathcal{F}_i$ est une tribu, $\forall i \in I, X \setminus A \in \mathcal{F}_i$, d'où $X \setminus A \in \bigcap_{i \in I} \mathcal{F}_i$.
   - Si $(A_n)$ est une suite dans l'intersection, elle est dans chaque $\mathcal{F}_i$, donc l'union est dans chaque $\mathcal{F}_i$, donc l'union est dans l'intersection.
3. Ainsi $\sigma(\mathcal{C})$ est bien une tribu contenant $\mathcal{C}$, et c'est manifestement la plus petite pour l'inclusion.

\textbf{Théorème :} Sur $\mathbb{R}$, la tribu borélienne $\mathcal{B}(\mathbb{R})$ est engendrée par la classe des intervalles semi-ouverts $\mathcal{C} = \{ ]-\infty, x] \mid x \in \mathbb{R} \}$.
\textit{Preuve détaillée :}
Notons $\mathcal{T}$ l'ensemble des ouverts de $\mathbb{R}$.
- Montrons $\sigma(\mathcal{C}) \subset \mathcal{B}(\mathbb{R})$ : chaque intervalle $]-\infty, x]$ est un fermé, donc son complémentaire $]x, +\infty[$ est un ouvert. Ainsi, $]-\infty, x]$ appartient à $\mathcal{B}(\mathbb{R})$. Puisque $\sigma(\mathcal{C})$ est la plus petite tribu contenant $\mathcal{C}$, $\sigma(\mathcal{C}) \subset \mathcal{B}(\mathbb{R})$.
- Montrons $\mathcal{B}(\mathbb{R}) \subset \sigma(\mathcal{C})$ :
  Pour tout $a < b$, on peut écrire l'intervalle semi-ouvert $]a, b] = ]-\infty, b] \setminus ]-\infty, a] = ]-\infty, b] \cap (]-\infty, a])^c$. Donc $]a, b] \in \sigma(\mathcal{C})$.
  L'intervalle ouvert $]a, b[$ s'écrit $]a, b[ = \bigcup_{n \geq 1} ]a, b - \frac{1}{n}]$. L'union étant dénombrable, $]a, b[ \in \sigma(\mathcal{C})$.
  Puisque tout ouvert de $\mathbb{R}$ s'écrit comme une union dénombrable d'intervalles ouverts, tout ouvert est dans $\sigma(\mathcal{C})$. La tribu de Borel étant engendrée par les ouverts, $\mathcal{B}(\mathbb{R}) \subset \sigma(\mathcal{C})$.
L'égalité est donc démontrée.

\section{Applications en Théorie des Probabilités et IA}

Dans la formalisation axiomatique de Kolmogorov (1933), un espace de probabilité est un triplet $(\Omega, \mathcal{F}, \mathbb{P})$.
- $\Omega$ est l'univers (les réalisations possibles).
- $\mathcal{F}$ est la tribu des événements (les ensembles auxquels on sait affecter une probabilité).
- $\mathbb{P} : \mathcal{F} \to [0, 1]$ est la mesure de probabilité.

En IA et en apprentissage automatique, les variables aléatoires (comme les poids d'un réseau ou les variables latentes dans un VAE) sont des fonctions mesurables.
Définition : Une fonction $X : \Omega \to E$ entre deux espaces mesurables $(\Omega, \mathcal{F})$ et $(E, \mathcal{E})$ est dite \textbf{mesurable} si pour tout $B \in \mathcal{E}$, on a $X^{-1}(B) \in \mathcal{F}$.
Cela garantit que l'on peut toujours demander la probabilité de l'événement $\{ X \in B \}$ !
Les probabilités modernes sur des espaces complexes (espaces de fonctions, graphes) requièrent l'utilisation méticuleuse des tribus cylindriques et de la tribu borélienne pour assurer une assise rigoureuse.


\part{Exercices d'Application et de Concours}
\subsection*{Exercice 1 : La tribu grossière et discrète \quad $\bigstar\star\star\star\star$}


\textbf{Énoncé :}
Soit $X$ un ensemble. Démontrer que l'intersection de toutes les tribus sur $X$ est la tribu grossière $\{\emptyset, X\}$, et que l'union de toutes les tribus n'est pas nécessairement une tribu, mais que la tribu des parties $\mathcal{P}(X)$ est la plus grande.

\textbf{Correction exhaustive pas-à-pas :}
1. Soit $\mathfrak{F}$ l'ensemble de toutes les tribus sur $X$. Par définition, pour toute tribu $\mathcal{F} \in \mathfrak{F}$, on a $\emptyset \in \mathcal{F}$ et $X \in \mathcal{F}$.
2. Par conséquent, $\{\emptyset, X\} \subset \bigcap_{\mathcal{F} \in \mathfrak{F}} \mathcal{F}$.
3. De plus, $\mathcal{F}_0 = \{\emptyset, X\}$ est elle-même une tribu sur $X$, donc elle appartient à $\mathfrak{F}$. L'intersection est donc contenue dans $\mathcal{F}_0$.
4. Ainsi, $\bigcap_{\mathcal{F} \in \mathfrak{F}} \mathcal{F} = \{\emptyset, X\}$.
5. Pour l'union : soient $\mathcal{F}_1 = \{\emptyset, \{1\}, \{2, 3\}, \{1, 2, 3\}\}$ et $\mathcal{F}_2 = \{\emptyset, \{2\}, \{1, 3\}, \{1, 2, 3\}\}$ sur $X=\{1, 2, 3\}$. $\mathcal{F}_1 \cup \mathcal{F}_2$ contient $\{1\}$ et $\{2\}$, mais pas leur union $\{1, 2\}$. Donc l'union de tribus n'est pas une tribu.



\subsection*{Exercice 2 : Tribu engendrée par un sous-ensemble \quad $\bigstar\bigstar\star\star\star$}


\textbf{Énoncé :}
Déterminer explicitement la tribu $\sigma(\{A\})$ engendrée par une partie $A \subset X$.

\textbf{Correction exhaustive pas-à-pas :}
1. Une tribu contenant $A$ doit contenir $A$.
2. Par stabilité au complémentaire, elle doit contenir $A^c = X \setminus A$.
3. Elle doit contenir l'ensemble vide $\emptyset$ et l'espace entier $X$.
4. Posons $\mathcal{T} = \{\emptyset, A, A^c, X\}$.
5. Vérifions que $\mathcal{T}$ est une tribu :
   - $X \in \mathcal{T}$.
   - Stabilité par complémentaire : $A^c \in \mathcal{T}, (A^c)^c = A \in \mathcal{T}, \emptyset^c = X \in \mathcal{T}, X^c = \emptyset \in \mathcal{T}$.
   - Stabilité par union : $A \cup A^c = X$, $A \cup \emptyset = A$, etc. Toutes les unions possibles d'éléments de $\mathcal{T}$ restent dans $\mathcal{T}$.
6. $\mathcal{T}$ est donc une tribu contenant $A$. C'est manifestement la plus petite, car toute tribu contenant $A$ doit contenir ces quatre éléments.
7. Ainsi, $\sigma(\{A\}) = \{\emptyset, A, A^c, X\}$.



\subsection*{Exercice 3 : Tribu co-dénombrable \quad $\bigstar\bigstar\star\star\star$}


\textbf{Énoncé :}
Soit $X$ un ensemble non dénombrable. Montrer que $\mathcal{F} = \{A \subset X \mid A \text{ est dénombrable ou } A^c \text{ est dénombrable}\}$ est une tribu sur $X$.

\textbf{Correction exhaustive pas-à-pas :}
1. L'ensemble vide $\emptyset$ est fini (donc dénombrable), donc $\emptyset \in \mathcal{F}$, et son complémentaire $X \in \mathcal{F}$.
2. Stabilité par complémentaire : Si $A \in \mathcal{F}$, alors par définition, soit $A$ est dénombrable (donc le complémentaire de $A^c$ est dénombrable, d'où $A^c \in \mathcal{F}$), soit $A^c$ est dénombrable (donc $A^c \in \mathcal{F}$).
3. Stabilité par union dénombrable : Soit $(A_n)_{n \in \mathbb{N}}$ une suite d'éléments de $\mathcal{F}$.
   - Cas 1 : Pour tout $n$, $A_n$ est dénombrable. Une réunion dénombrable d'ensembles dénombrables est dénombrable, donc $\bigcup A_n$ est dénombrable, d'où $\bigcup A_n \in \mathcal{F}$.
   - Cas 2 : Il existe au moins un entier $k$ tel que $A_k^c$ est dénombrable.
     Alors $(\bigcup A_n)^c = \bigcap A_n^c \subset A_k^c$.
     Puisque $A_k^c$ est dénombrable, tout sous-ensemble de $A_k^c$ est dénombrable.
     Donc $(\bigcup A_n)^c$ est dénombrable, ce qui implique $\bigcup A_n \in \mathcal{F}$.
4. Conclusion : $\mathcal{F}$ est bien une tribu.



\subsection*{Exercice 4 : Propriétés de l'opérateur $\sigma$ \quad $\bigstar\bigstar\bigstar\star\star$}


\textbf{Énoncé :}
Soient $\mathcal{C}_1$ et $\mathcal{C}_2$ deux familles de parties de $X$. Montrer que :
1. Si $\mathcal{C}_1 \subset \mathcal{C}_2$, alors $\sigma(\mathcal{C}_1) \subset \sigma(\mathcal{C}_2)$.
2. $\sigma(\sigma(\mathcal{C}_1)) = \sigma(\mathcal{C}_1)$.

\textbf{Correction exhaustive pas-à-pas :}
1. $\sigma(\mathcal{C}_2)$ est une tribu qui contient $\mathcal{C}_2$. Puisque $\mathcal{C}_1 \subset \mathcal{C}_2$, on a $\mathcal{C}_1 \subset \sigma(\mathcal{C}_2)$.
   Or $\sigma(\mathcal{C}_1)$ est l'intersection de toutes les tribus contenant $\mathcal{C}_1$. Comme $\sigma(\mathcal{C}_2)$ est l'une de ces tribus, on a nécessairement $\sigma(\mathcal{C}_1) \subset \sigma(\mathcal{C}_2)$.
2. Posons $\mathcal{T} = \sigma(\mathcal{C}_1)$. $\mathcal{T}$ est une tribu. L'opérateur $\sigma$ associe à une famille la plus petite tribu la contenant. Puisque $\mathcal{T}$ est déjà une tribu, la plus petite tribu contenant $\mathcal{T}$ est $\mathcal{T}$ elle-même.
   Ainsi, $\sigma(\sigma(\mathcal{C}_1)) = \sigma(\mathcal{T}) = \mathcal{T} = \sigma(\mathcal{C}_1)$.



\subsection*{Exercice 5 : Tribu trace \quad $\bigstar\bigstar\bigstar\star\star$}


\textbf{Énoncé :}
Soit $(X, \mathcal{F})$ un espace mesurable et $Y \subset X$ un sous-ensemble (pas nécessairement dans $\mathcal{F}$). Démontrer que $\mathcal{F}_Y = \{A \cap Y \mid A \in \mathcal{F}\}$ est une tribu sur $Y$, appelée tribu trace.

\textbf{Correction exhaustive pas-à-pas :}
1. L'ensemble entier pour $Y$ est $Y$. Or $X \in \mathcal{F}$ et $X \cap Y = Y$, donc $Y \in \mathcal{F}_Y$.
2. Stabilité par complémentaire relatif : Soit $B \in \mathcal{F}_Y$. Il existe $A \in \mathcal{F}$ tel que $B = A \cap Y$.
   Le complémentaire de $B$ dans $Y$ est $Y \setminus B = Y \setminus (A \cap Y) = Y \setminus A = (X \setminus A) \cap Y$.
   Puisque $\mathcal{F}$ est une tribu, $X \setminus A \in \mathcal{F}$. Donc $(X \setminus A) \cap Y \in \mathcal{F}_Y$.
3. Stabilité par union dénombrable : Soit $(B_n)_{n \in \mathbb{N}}$ une suite dans $\mathcal{F}_Y$. Pour tout $n$, il existe $A_n \in \mathcal{F}$ tel que $B_n = A_n \cap Y$.
   $\bigcup_{n \in \mathbb{N}} B_n = \bigcup_{n \in \mathbb{N}} (A_n \cap Y) = \left( \bigcup_{n \in \mathbb{N}} A_n \right) \cap Y$.
   Puisque $\mathcal{F}$ est une tribu, $\bigcup A_n \in \mathcal{F}$. Ainsi, l'union des $B_n$ appartient à $\mathcal{F}_Y$.
4. $\mathcal{F}_Y$ est bien une tribu sur $Y$.



\subsection*{Exercice 6 : Atomes d'une tribu \quad $\bigstar\bigstar\bigstar\bigstar\star$}


\textbf{Énoncé :}
Soit $X$ un ensemble et $\mathcal{F}$ une tribu engendrée par une partition dénombrable $\{E_n\}_{n \in \mathbb{N}}$ de $X$ ($E_n \neq \emptyset$, disjoints, réunion vaut $X$). Montrer que tout élément de $\mathcal{F}$ est une union (éventuellement vide ou dénombrable) d'éléments $E_n$.

\textbf{Correction exhaustive pas-à-pas :}
1. Soit $\mathcal{T}$ l'ensemble des unions, finies, dénombrables ou vides, des éléments $E_n$.
2. Montrons que $\mathcal{T}$ est une tribu.
   - $\emptyset \in \mathcal{T}$ (union vide). $X \in \mathcal{T}$ car $X = \bigcup E_n$.
   - Stabilité par complémentaire : Soit $A \in \mathcal{T}$, $A = \bigcup_{n \in J} E_n$ où $J \subset \mathbb{N}$. Le complémentaire de $A$ est $X \setminus A = \bigcup_{n \notin J} E_n$ (puisque les $E_n$ forment une partition). Cette union appartient bien à $\mathcal{T}$.
   - Stabilité par union dénombrable : Une union dénombrable d'unions dénombrables d'éléments $E_n$ reste une union dénombrable d'éléments $E_n$.
3. $\mathcal{T}$ est donc une tribu. De plus, chaque $E_k \in \mathcal{T}$, donc $\mathcal{T}$ contient la partition.
4. Comme $\mathcal{F}$ est la plus petite tribu contenant la partition, on a $\mathcal{F} \subset \mathcal{T}$.
5. D'autre part, toute tribu contenant les $E_n$ doit contenir leurs unions dénombrables, donc $\mathcal{T} \subset \mathcal{F}$.
6. Par conséquent, $\mathcal{F} = \mathcal{T}$. Les $E_n$ sont appelés les atomes de cette tribu.



\subsection*{Exercice 7 : Générateurs de la tribu borélienne \quad $\bigstar\bigstar\bigstar\bigstar\star$}


\textbf{Énoncé :}
Démontrer que la tribu de Borel sur $\mathbb{R}$, $\mathcal{B}(\mathbb{R})$, est engendrée par les intervalles ouverts rationnels $]p, q[$ avec $p, q \in \mathbb{Q}$.

\textbf{Correction exhaustive pas-à-pas :}
1. Notons $\mathcal{D} = \{ ]p, q[ \mid p, q \in \mathbb{Q}, p < q \}$. Puisque chaque intervalle $]p, q[$ est un ouvert, il appartient à $\mathcal{B}(\mathbb{R})$. Donc $\sigma(\mathcal{D}) \subset \mathcal{B}(\mathbb{R})$.
2. Pour l'inclusion inverse, il suffit de montrer que tout ouvert de $\mathbb{R}$ appartient à $\sigma(\mathcal{D})$.
3. Tout ouvert de $\mathbb{R}$ peut s'écrire comme une réunion au plus dénombrable d'intervalles ouverts disjoints : $U = \bigcup_{i} ]a_i, b_i[$.
4. Soit un intervalle quelconque $]a, b[$. Par la densité de $\mathbb{Q}$ dans $\mathbb{R}$, on peut trouver deux suites de rationnels $(p_n)$ et $(q_n)$ telles que $p_n \downarrow a$ et $q_n \uparrow b$ avec $p_n < q_n$.
5. On a alors $]a, b[ = \bigcup_{n} ]p_n, q_n[$. Chaque $]p_n, q_n[ \in \mathcal{D} \subset \sigma(\mathcal{D})$, donc l'union dénombrable $]a, b[ \in \sigma(\mathcal{D})$.
6. Ainsi, $U \in \sigma(\mathcal{D})$. Puisque $\mathcal{B}(\mathbb{R})$ est générée par les ouverts, $\mathcal{B}(\mathbb{R}) \subset \sigma(\mathcal{D})$.
7. En conclusion, $\mathcal{B}(\mathbb{R}) = \sigma(\mathcal{D})$.



\subsection*{Exercice 8 : Image réciproque de tribu \quad $\bigstar\bigstar\bigstar\bigstar\bigstar$}


\textbf{Énoncé :}
Soit $f : X \to Y$ une application et $\mathcal{F}_Y$ une tribu sur $Y$. Montrer que l'ensemble $f^{-1}(\mathcal{F}_Y) = \{f^{-1}(B) \mid B \in \mathcal{F}_Y\}$ est une tribu sur $X$.

\textbf{Correction exhaustive pas-à-pas :}
1. L'espace complet : $Y \in \mathcal{F}_Y$. L'image réciproque est $f^{-1}(Y) = X$. Donc $X \in f^{-1}(\mathcal{F}_Y)$.
2. Stabilité par complémentaire : Soit $A \in f^{-1}(\mathcal{F}_Y)$. Il existe $B \in \mathcal{F}_Y$ tel que $A = f^{-1}(B)$.
   Le complémentaire est $X \setminus A = X \setminus f^{-1}(B) = f^{-1}(Y \setminus B)$.
   Puisque $\mathcal{F}_Y$ est une tribu, $Y \setminus B \in \mathcal{F}_Y$. Donc $X \setminus A \in f^{-1}(\mathcal{F}_Y)$.
3. Stabilité par union dénombrable : Soit $(A_n)_{n \in \mathbb{N}}$ une suite dans $f^{-1}(\mathcal{F}_Y)$. Il existe une suite $(B_n)$ dans $\mathcal{F}_Y$ telle que $A_n = f^{-1}(B_n)$.
   $\bigcup A_n = \bigcup f^{-1}(B_n) = f^{-1}(\bigcup B_n)$.
   Puisque $\mathcal{F}_Y$ est une tribu, $\bigcup B_n \in \mathcal{F}_Y$. Ainsi, $\bigcup A_n \in f^{-1}(\mathcal{F}_Y)$.
4. Conclusion : $f^{-1}(\mathcal{F}_Y)$ est bien une tribu sur $X$.



\subsection*{Exercice 9 : Tribu produit (Introduction) \quad $\bigstar\bigstar\bigstar\bigstar\bigstar$}


\textbf{Énoncé :}
Soient $(X, \mathcal{A})$ et $(Y, \mathcal{B})$ deux espaces mesurables. On note $\mathcal{R} = \{A \times B \mid A \in \mathcal{A}, B \in \mathcal{B}\}$ la famille des pavés mesurables. Montrer que $\mathcal{R}$ n'est en général pas une algèbre (donc pas une tribu).

\textbf{Correction exhaustive pas-à-pas :}
1. Soient $A \times B \in \mathcal{R}$. Son complémentaire dans $X \times Y$ est :
   $(X \times Y) \setminus (A \times B) = ((X \setminus A) \times Y) \cup (X \times (Y \setminus B))$.
2. En général, la réunion de deux pavés n'est pas un pavé. Prenons $X = Y = \mathbb{R}$, avec $\mathcal{A} = \mathcal{B} = \mathcal{B}(\mathbb{R})$.
   Soit $A = [0, 1]$ et $B = [0, 1]$. $A \times B$ est le carré unité.
3. Son complémentaire est la région plane privée du carré. Cette région ne peut pas s'écrire sous la forme d'un unique produit cartésien $U \times V$.
4. Puisque le complémentaire d'un élément de $\mathcal{R}$ n'est pas dans $\mathcal{R}$, la famille $\mathcal{R}$ n'est pas stable par passage au complémentaire.
5. Elle n'est donc pas une algèbre, et \textit{a fortiori} pas une tribu.



\subsection*{Exercice 10 : Lemme des classes monotones (Cas simplifé) \quad $\bigstar\bigstar\bigstar\bigstar\bigstar$}


\textbf{Énoncé :}
Une classe monotone $\mathcal{M}$ est stable par limites croissantes dénombrables et limites décroissantes dénombrables. Démontrer que si une algèbre $\mathcal{A}$ est une classe monotone, alors c'est une tribu.

\textbf{Correction exhaustive pas-à-pas :}
1. On sait que $\mathcal{A}$ est une algèbre : elle contient $X$, est stable par complémentaire et par union finie.
2. Pour montrer que c'est une tribu, il suffit de prouver la stabilité par union dénombrable.
3. Soit $(A_n)_{n \in \mathbb{N}}$ une suite d'éléments de $\mathcal{A}$.
4. Posons $B_N = \bigcup_{n=0}^N A_n$. Puisque $\mathcal{A}$ est une algèbre (stable par union finie), $B_N \in \mathcal{A}$ pour tout $N$.
5. La suite $(B_N)$ est croissante au sens de l'inclusion : $B_0 \subset B_1 \subset B_2 \dots$
6. La limite croissante de cette suite est exactement $\bigcup_{n=0}^{\infty} A_n$.
7. Puisque $\mathcal{A}$ est une classe monotone, la limite croissante d'une suite d'éléments de $\mathcal{A}$ appartient à $\mathcal{A}$.
8. Par conséquent, $\bigcup_{n=0}^{\infty} A_n \in \mathcal{A}$.
9. L'algèbre $\mathcal{A}$ est donc stable par union dénombrable, ce qui prouve que c'est une tribu.





\part{Travaux Pratiques et Simulations Algorithmiques}
\subsection*{TP 1 : Implémentation ensembliste naïve d'une tribu finie}

\begin{lstlisting}

# Représentation d'une tribu générée par une partition finie
def generate_algebra(partition):
    # La partition est une liste de sets disjoints
    universe = set().union(*partition)

    # Générer toutes les unions possibles (2^n éléments)
    # L'union de combinaisons de la partition garantit la structure de tribu
    n = len(partition)
    algebra = []

    for i in range(1 << n):
        element = set()
        for j in range(n):
            if (i & (1 << j)):
                element = element.union(partition[j])
        algebra.append(frozenset(element))

    return set(algebra)

# Validation
P = [{"a", "b"}, {"c"}, {"d", "e"}]
tribu = generate_algebra(P)

print("Taille de la tribu :", len(tribu)) # 2^3 = 8
assert len(tribu) == 8

# Vérification de la stabilité par complémentaire
universe = frozenset().union(*P)
for A in tribu:
    complement = universe - A
    assert complement in tribu
print("Stabilité par complémentaire validée.")

\end{lstlisting}


\subsection*{TP 2 : Génération de la tribu borélienne sur une grille discrète}

\begin{lstlisting}

# Approximation discrète de la tribu borélienne 1D
def generate_topology_on_grid(grid_size):
    # Les ouverts de base sont les 'intervalles ouverts' i.e. segments continus sans bords
    # Sur une grille discrète, on simule l'ouverture stricte.
    base_open_sets = []
    for start in range(grid_size):
        for end in range(start + 2, grid_size + 1):
            # Un segment ouvert nécessite au moins 1 point strict entre start et end
            base_open_sets.append(frozenset(range(start + 1, end - 1)))

    return base_open_sets

def closure_union_complement(generator_sets, universe):
    current = set(generator_sets)
    current.add(frozenset(universe))
    current.add(frozenset())

    added = True
    while added:
        added = False
        new_sets = set()

        # Complémentaire
        for A in current:
            Ac = frozenset(universe - A)
            if Ac not in current and Ac not in new_sets:
                new_sets.add(Ac)

        # Unions paires
        current_list = list(current)
        for i in range(len(current_list)):
            for j in range(i, len(current_list)):
                union_set = current_list[i].union(current_list[j])
                if union_set not in current and union_set not in new_sets:
                    new_sets.add(union_set)

        if new_sets:
            current.update(new_sets)
            added = True

    return current

grid_size = 5
universe = set(range(grid_size))
base = generate_topology_on_grid(grid_size)
borel = closure_union_complement(base, universe)

print("Taille de la tribu borélienne discrète :", len(borel))
# Dans ce cas fini, Borel = Parties(X) car les singletons sont des fermés, donc leur union donne tout.
assert len(borel) == 2**grid_size

\end{lstlisting}


\subsection*{TP 3 : Validation de la stabilité d'une tribu (Automatisée)}

\begin{lstlisting}

def is_algebra(family, universe):
    # family est un set de frozensets

    # 1. Vérifier si l'univers et l'ensemble vide sont présents
    if frozenset(universe) not in family or frozenset() not in family:
        return False

    # 2. Vérifier la stabilité par complémentaire
    for A in family:
        complement = frozenset(universe - A)
        if complement not in family:
            return False

    # 3. Vérifier la stabilité par union finie
    for A in family:
        for B in family:
            union_set = A.union(B)
            if union_set not in family:
                return False

    return True

# Test 1 : L'ensemble des parties
U1 = {1, 2}
F1 = {frozenset(), frozenset({1}), frozenset({2}), frozenset({1, 2})}
assert is_algebra(F1, U1) == True

# Test 2 : Manque un complémentaire
F2 = {frozenset(), frozenset({1}), frozenset({1, 2})}
assert is_algebra(F2, U1) == False

# Test 3 : Manque l'union
U3 = {1, 2, 3}
F3 = {frozenset(), frozenset({1}), frozenset({2, 3}), frozenset({2}), frozenset({1, 3}), frozenset(U3)}
assert is_algebra(F3, U3) == False # {1} U {2} = {1, 2} n'est pas dans F3

print("Vérificateur d'algèbre opérationnel.")

\end{lstlisting}


\subsection*{TP 4 : Approximation de la tribu trace}

\begin{lstlisting}

def trace_algebra(family, subset):
    # Génère la tribu trace sur le sous-ensemble 'subset'
    trace = set()
    for A in family:
        intersection = frozenset(A.intersection(subset))
        trace.add(intersection)
    return trace

# Définissons un univers et une tribu valide (générée par une partition)
universe = {1, 2, 3, 4, 5, 6}
partition = [{"1", "2"}, {"3", "4"}, {"5", "6"}]
# Remplace les strings par des int pour matcher l'univers
P = [{1, 2}, {3, 4}, {5, 6}]

F = set()
F.add(frozenset())
F.add(frozenset(universe))
F.add(frozenset({1, 2}))
F.add(frozenset({3, 4}))
F.add(frozenset({5, 6}))
F.add(frozenset({1, 2, 3, 4}))
F.add(frozenset({1, 2, 5, 6}))
F.add(frozenset({3, 4, 5, 6}))

subset = {2, 3, 6}
trace = trace_algebra(F, subset)

print("Éléments de la tribu trace :")
for el in trace:
    print(set(el))

# Validation : la trace doit elle-même être une tribu sur 'subset'
def is_algebra_trace(family, local_universe):
    if frozenset(local_universe) not in family: return False
    for A in family:
        if frozenset(local_universe - A) not in family: return False
        for B in family:
            if A.union(B) not in family: return False
    return True

assert is_algebra_trace(trace, subset) == True
print("Propriété de la tribu trace validée.")

\end{lstlisting}


\subsection*{TP 5 : Image réciproque d'une algèbre}

\begin{lstlisting}

def inverse_image_algebra(mapping, Y_family, domain):
    # mapping est un dictionnaire représentant une fonction de domain dans codomain
    inv_family = set()

    for B in Y_family:
        preimage = set()
        for x in domain:
            if mapping.get(x) in B:
                preimage.add(x)
        inv_family.add(frozenset(preimage))

    return inv_family

# Soit f : X -> Y
X = {"a", "b", "c", "d"}
Y = {1, 2, 3}
f = {"a": 1, "b": 1, "c": 2, "d": 3}

# Soit F_Y une tribu sur Y (la tribu des parties)
F_Y = {
    frozenset(), frozenset({1}), frozenset({2}), frozenset({3}),
    frozenset({1, 2}), frozenset({1, 3}), frozenset({2, 3}), frozenset(Y)
}

inv_F = inverse_image_algebra(f, F_Y, X)

# Validation que l'image réciproque est une tribu sur X
def validate_algebra(family, univ):
    if frozenset(univ) not in family: return False
    for A in family:
        if frozenset(univ - A) not in family: return False
    for A in family:
        for B in family:
            if A.union(B) not in family: return False
    return True

assert validate_algebra(inv_F, X) == True
print("L'image réciproque de la tribu a conservé sa structure.")

\end{lstlisting}




\end{document}
